\input{../../../stock/en-tete_v6.tex}

\newcommand{\titrechapitre}{Logique des prédicats}
\newcommand{\numerochapitre}{3}

\renewcommand{\headrulewidth}{0.4pt}
\renewcommand{\footrulewidth}{0.4pt}
\fancyfoot[L]{Notes de Raphaël JONTEF}
\fancyfoot[C]{\thepage}
\fancyfoot[R]{Enseirb-Matmeca}
\fancyhead[L]{Cours de Mathématiques}
\fancyhead[R]{\titrechapitre}

\begin{document}
\input{../../../stock/commands.tex}

% Page de titre custom
\setlength{\headheight}{15pt}
\begin{titlepage}
    \centering
    \vspace*{3cm}
    {\huge Chapitre \numerochapitre\par}
    {\Huge \bfseries\titrechapitre\par}
    \vspace{2cm}
    {\Large Notes rédigées par Raphaël JONTEF\par}
{\Large Avec l'aide de Github Copilot\par}
    \vspace{0.5cm}
    {\normalsize Cours enseigné par Frédéric HERBRETEAU\par}
    \vspace{4cm}
    {\small Enseirb-Matmeca\par}
    {\small filière informatique\par}
    \vfill
    {\large \today\par}
\end{titlepage}

% plan du cours
\tableofcontents
\newpage


% -----------Corps du document--------------

\section{Quantificateurs}
\begin{definition}{}{prédicat}
    Un \notion{prédicat} est une fonction qui prend un ou plusieurs arguments et retourne une valeur de vérité (vrai ou faux).
\end{definition}

\begin{exemple}{}{}
    le prédicat $P(x)$ peut être défini comme ``$x$ est un nombre pair''. Si l'on considère le domaine des entiers, alors $P(2)$ est vrai, tandis que $P(3)$ est faux.
\end{exemple}

\begin{definition}{}{quantificateur universel}
    Le \notion{quantificateur universel} est noté $\forall$ et signifie "pour tout". Par exemple, l'expression $\forall x P(x)$ signifie que le prédicat $P(x)$ est vrai pour tous les éléments $x$ dans le domaine de discours.
\end{definition}

\begin{definition}{}{quantificateur existentiel}
    Le \notion{quantificateur existentiel} est noté $\exists$ et signifie "il existe". Par exemple, l'expression $\exists x P(x)$ signifie qu'il existe au moins un élément $x$ dans le domaine de discours pour lequel le prédicat $P(x)$ est vrai.
\end{definition}

\begin{remarque}{}{}
    Dans un prédicat, l'ordre des quantificateurs est crucial. Par exemple, les expressions $\forall x \exists y P(x, y)$ et $\exists y \forall x P(x, y)$ peuvent avoir des significations très différentes~:
    \begin{center}
        "Pour tout enfant il existe une maman qui l'aime."
    \end{center}
    \begin{center}
        "Il existe une maman qui aime tous les enfants."
    \end{center}
    Ces deux prédicats ne sont visiblement pas équivalents : l'un est utopique, l'autre est déroutant.
\end{remarque}

\begin{remarque}{}{}
    Par convention, une assertion universelle sur un domaine vide est considérée comme vraie, tandis qu'une assertion existentielle sur un domaine vide est considérée comme fausse.
\end{remarque}


\begin{definition}{}{logique du premier ordre, langage du premier ordre}
    La \notion{logique du premier ordre} (ou logique des prédicats) est une extension de la logique propositionnelle, dans laquelle on travaille sur un \notion{langage du premier ordre}, défini par~:
    \begin{itemize}
        \item des symboles de constantes (comme $0$, $\varnothing$, $\pi$,...)
        \item des symboles de fonctions, comme $+ (5,3)$
        \item des symboles de relations, comme $=$, $\neq$, $\geq$, $>$.
    \end{itemize}
    où l'on introduit aussi des quantificateurs. Les constantes et images des symboles de fonctions sont des éléments d'un ensemble d'étude, appelé \notion{domaine}, qui dépend du langage du premier ordre.
\end{definition}

\begin{definition}{}{occurrence libre, occurrence liée}
    Dans une formule de logique du premier ordre, une \notion{occurrence libre} d'une variable est une occurrence qui n'est pas introduite dans un quantificateur. Une \notion{occurrence liée} est une occurrence qui est sous la portée d'un quantificateur.
\end{definition}

\begin{definition}{}{interprétation}
    Une \notion{interprétation} pour un langage du premier assigne des significations aux symboles de constantes, fonctions et relations dans un domaine donné~:
    \begin{itemize}
        \item Chaque symbole de constante est associé à un élément spécifique du domaine.
        \item Chaque symbole de fonction est associé à une fonction définie sur le domaine.
        \item Chaque symbole de relation est associé à une relation définie sur le domaine.
\end{definition}

\begin{remarque}{}{}
    On peut voir une interprétation comme un analogue des valuations en logique propositionnelle, mais adapté à la logique du premier ordre.
\end{remarque}

\begin{definition}{}{forme normale négative}
    Une formule est en \notion{forme normale négative (FNN)} si~:
    \begin{itemize}
        \item les seuls connecteurs logiques qu'elle utilise sont $\land$, $\lor$ et des quantificateurs.
        \item les négations n'apparaissent que directement devant des variables propositionnelles (formant des littéraux).
    \end{itemize} 
\end{definition}

\begin{definition}{}{satisfiabilité}
    Une formule est \notion{satisfiable} s'il existe une interprétation qui la rend vraie.
\end{definition}

\begin{remarque}{}{}
    Même si on a le choix sur l'interprétation, il existe des prédicats qui n'admettent aucune interprétation les rendant vrais. Par exemple, le prédicat $\forall x \Big(P(x) \land \lnegP(x)\Big)$ est insatisfiable.
    De plus, la satisfiabilité en logique du premier ordre est  un problème algorithmique indécidable.
\end{remarque}


\begin{definition}{}{validité}
    Une formule est \notion{valide} si elle est vraie dans toutes les interprétations possibles.
\end{definition}

\begin{exemple}{}{}
    Montrons que le prédicat suivant est satisfiable~:
    $$\forall x P(x) \implies \exists x P(x)$$
    Considérons une interprétation satisfaisant $\forall x P(x)$. Pour une telle interprétation, il existe au moins un élément $a$ dans le domaine tel que $P(a)$ est vrai. Donc, cette interprétation satisfait aussi $\exists x P(x)$.\\
    Le prédicat est donc valide.
\end{exemple}


\end{document}